% NNGraph DSL — گزارش فارسی
% کامپایل با: xelatex -shell-escape report-fa.tex
% (دو بار کامپایل برای فهرست مطالب و مراجع)

\documentclass{nngraph-fa}

% xepersian باید آخرین بسته بارگذاری شود
\usepackage{xepersian}

% تنظیم فونت — بعد از بارگذاری xepersian
\settextfont[
    BoldFont       = Vazirmatn Bold,
    ItalicFont     = Vazirmatn,
    BoldItalicFont = Vazirmatn Bold,
]{Vazirmatn}
\setdigitfont{Vazirmatn}
\setmonofont{Liberation Mono}

% ----------------------------------------------------------------
% اطلاعات سند
\title{
    \textbf{کامپایلر \en{NNGraph DSL}} \\[0.5em]
    \large یک زبان دامنه‌محور برای توصیف معماری شبکه‌های عصبی \\
    و تولید خودکار کد \en{PyTorch}
}
\author{
    \large گزارش فنی پروژه
}
\date{\today}

% ----------------------------------------------------------------
\begin{document}

\maketitle
\tableofcontents
\newpage

% ================================================================
\chapter*{چکیده}
\markboth{چکیده}{چکیده}
\addcontentsline{toc}{chapter}{چکیده}

\en{NNGraph DSL} یک زبان دامنه‌محور (\en{Domain-Specific Language}) است که برای توصیف
معماری شبکه‌های عصبی به صورت گراف جهت‌دار طراحی شده است. فایل‌های \code{.nng} را به
عنوان ورودی دریافت کرده و کد \en{PyTorch} اجرایی تولید می‌کند که یک زیرکلاس
\code{nn.Module} کامل با متدهای \code{\_\_init\_\_} و \code{forward} است.

کامپایلر با \en{ANTLR4} نسخه ۴.۱۳.۲ ساخته شده و از الگوی \en{Listener} برای تحلیل
معنایی و الگوریتم مرتب‌سازی توپولوژیک \en{Kahn} هم در تحلیل‌گر معنایی (تشخیص دور)
و هم در مولد کد (ترتیب اجرا) استفاده می‌کند.

مجموعه آزمایش شامل ۵۵ تست در چهار فایل است. چهار معماری نمونه---\en{MLP}،
\en{ResBlock}، \en{InceptionBlock} و \en{TransformerEncoder}---با موفقیت کامپایل
شده و \en{forward pass} آن‌ها با شکل خروجی صحیح اجرا شده است.

% ================================================================
\chapter{مقدمه}

\section{انگیزه}

نوشتن زیرکلاس‌های \code{nn.Module} در \en{PyTorch} به صورت دستی برای معماری‌های
پیچیده، تکراری و مستعد خطاست. برای یک شبکه با ده لایه، برنامه‌نویس باید:

\begin{itemize}
    \item در \code{\_\_init\_\_} هر لایه را جداگانه با پارامترهای صحیح تعریف کند
    \item در \code{forward} ترتیب اجرا را به صورت دستی مدیریت کند
    \item در توپولوژی‌های شاخه‌دار، متغیرهای میانی را نام‌گذاری کند
    \item مطمئن شود هیچ تنسور مشترکی پیش از استفاده بازنویسی نمی‌شود
\end{itemize}

\begin{exbox}{نمونه boilerplate دستی برای یک ResBlock}{resblock-manual}
برنامه‌نویس باید این ساختار را برای یک \en{ResBlock} ساده بنویسد:

\begin{latin}
\begin{verbatim}
class ResBlock(nn.Module):
    def __init__(self):
        super(ResBlock, self).__init__()
        self.conv1 = nn.Conv2d(in_channels=64, out_channels=64,
                               kernel_size=3, stride=1, padding=1)
        self.bn1   = nn.BatchNorm2d(num_features=64)
        self.relu1 = nn.ReLU()
        # ... و به همین ترتیب

    def forward(self, x):
        identity = x
        x = self.conv1(x)
        x = self.bn1(x)
        x = self.relu1(x)
        # ... مدیریت دستی اتصال باقیمانده
        x = x + identity
        return x
\end{verbatim}
\end{latin}

با \en{NNGraph DSL}، معماری یک‌بار به صورت گراف توصیف می‌شود و کامپایلر این کد را
تولید می‌کند.
\end{exbox}

\section{هدف}

هدف \en{NNGraph DSL} این است که فایل \code{.nng} تنها منبع حقیقت معماری باشد.
برنامه‌نویس گراف را توصیف می‌کند؛ کامپایلر تمام \en{boilerplate} را تولید می‌کند.

\section{دامنه کاربرد}

این ابزار برای توصیف معماری‌های ثابت (نه پویا) مناسب است---شبکه‌هایی که گراف
محاسباتی آن‌ها در زمان تعریف مدل مشخص است. پشتیبانی از کنترل جریان پویا
(\code{if}/\code{else} درون \code{forward}) در نسخه‌های آینده برنامه‌ریزی شده است.

% ================================================================
\chapter{کارهای مرتبط}

\section{چارچوب NADER}

\begin{dfnbox}{NADER}{nader}
\dfntxt{\en{NADER}} (\en{Neural Architecture Design via Representation}) چارچوبی است که
معماری شبکه‌های عصبی را به صورت گراف محاسباتی جهت‌دار (\en{DAG}) مدل می‌کند.
هر گره یک عملیات لایه است و هر یال جریان تنسور را نشان می‌دهد.
\end{dfnbox}

\en{NNGraph DSL} این مدل گراف را به عنوان پایه طراحی زبان به کار می‌گیرد.

\section{ANTLR4}

\begin{dfnbox}{ANTLR4}{antlr4}
\dfntxt{\en{ANTLR4}} (\en{ANother Tool for Language Recognition}) ابزاری است که از یک
فایل گرامر \code{.g4} به صورت خودکار \en{lexer}، \en{parser} و \en{listener/visitor}
تولید می‌کند. نسخه مورد استفاده: \en{ANTLR 4.13.2}.
\end{dfnbox}

\subsection{الگوی Listener در برابر Visitor}

\en{ANTLR4} دو الگو برای پیمایش درخت تجزیه ارائه می‌دهد:

\begin{itemize}
    \item \textbf{\en{Listener}}: \code{ParseTreeWalker} به صورت خودکار درخت را پیمایش
    می‌کند و متدهای \code{enter*}/\code{exit*} را فراخوانی می‌کند. هیچ مقدار بازگشتی
    نیاز نیست.
    \item \textbf{\en{Visitor}}: برنامه‌نویس پیمایش را کنترل می‌کند و هر متد یک مقدار
    برمی‌گرداند.
\end{itemize}

\begin{notebox}
این کامپایلر از الگوی \en{Listener} استفاده می‌کند. متدهای \code{exit*} در
\code{custom\_listener.py} حالت را در \code{self.nodes} و \code{self.edges}
انباشته می‌کنند---این رویکرد برای ساختن یک مدل گراف از یک درخت تجزیه
طبیعی‌تر است.
\end{notebox}

\section{پروژه EVM ANTLR}

این کامپایلر سبک و ساختار کد را از پروژه \en{EVM} موجود در
\code{PycharmProjects/Antlr} به ارث می‌برد. سه فایل از آن پروژه مستقیماً
استفاده می‌شوند:

\begin{itemize}
    \item \code{ast.py} --- کلاس \code{AST} و \code{TreeNode}
    \item \code{make\_ast\_subtree.py} --- ترکیب گره‌های درخت تجزیه
    \item \code{ast\_to\_networkx\_graph.py} --- تجسم اختیاری گراف
\end{itemize}

% ================================================================
\chapter{طراحی زبان}

\section{ساختار برنامه}

هر فایل \code{.nng} از سه بلوک تشکیل شده است. بلوک \code{config} اختیاری است:

\begin{latin}
\begin{verbatim}
model <Name> {
    input  <id> : tensor(<shape>)
    output <id>
}

graph {
    node <id> : <LayerType>(<params>)
    ...
    edge <src> -> <dst>
    edge <src> -> <dst> [label="<string>"]
    ...
}

config {
    batch_size = <int>
    device     = "cpu" | "cuda"
}
\end{verbatim}
\end{latin}

\begin{dfnbox}{بلوک model}{block-model}
بلوک \code{model} نام کلاس تولیدشده، شناسه گره ورودی و شکل تنسور ورودی، و
شناسه گره خروجی را مشخص می‌کند.
\end{dfnbox}

\begin{dfnbox}{بلوک graph}{block-graph}
بلوک \code{graph} تمام گره‌های لایه (\code{node}) و یال‌های جهت‌دار (\code{edge})
را تعریف می‌کند. ترتیب اعلان یال‌ها اهمیتی ندارد؛ کامپایلر ترتیب اجرا را از
ساختار گراف محاسبه می‌کند.
\end{dfnbox}

\begin{dfnbox}{بلوک config}{block-config}
بلوک \code{config} مقادیری را برای بلوک \code{\_\_main\_\_} کد تولیدشده تعیین
می‌کند. \code{batch\_size} پیش‌فرض ۱ و \code{device} پیش‌فرض \code{'cpu'} است.
\end{dfnbox}

\section{گرامر}

\begin{dfnbox}{آمار گرامر NNGraph.g4}{grammar-stats}
\begin{itemize}[noitemsep]
    \item فایل گرامر: \code{NNGraph.g4} (۵۶ خط)
    \item قوانین تجزیه‌گر: ۱۴ قانون (\en{snake\_case})
    \item توکن‌های واژه‌نگار: ۲۰ توکن (\en{ALL\_CAPS})
\end{itemize}
\end{dfnbox}

قوانین تجزیه‌گر:

\begin{latin}
\begin{verbatim}
program     : model_block graph_block config_block? EOF;
model_block : MODEL ID '{' input_decl output_decl '}';
input_decl  : INPUT ID ':' TENSOR shape_expr;
output_decl : OUTPUT ID;
graph_block : GRAPH '{' (node_decl | edge_decl)* '}';
node_decl   : NODE ID ':' layer_expr;
layer_expr  : ID '(' param_list? ')';
param_list  : param (',' param)*;
param       : ID '=' value;
edge_decl   : EDGE ID ARROW ID ('[' LABEL '=' STRING ']')?;
config_block: CONFIG '{' config_entry* '}';
config_entry: ID '=' value;
value       : FLOAT_LITERAL | INT_LITERAL | BOOL_LITERAL
            | STRING | NONE | shape_expr;
shape_expr  : '(' INT_LITERAL (',' INT_LITERAL)* ')';
\end{verbatim}
\end{latin}

\begin{notebox}
\textbf{نکته مهم طراحی:} \code{input\_decl} از \code{TENSOR shape\_expr} استفاده
می‌کند نه \code{TENSOR '(' shape\_expr ')'}، زیرا \code{shape\_expr} خود پرانتزها
را در بر می‌گیرد. این خطای رایج در نسخه‌های اولیه گرامر برطرف شد.
\end{notebox}

\section{سیستم انواع}

\begin{center}
\begin{tabular}{lll}
\toprule
\textbf{نوع در DSL} & \textbf{نوع Python} & \textbf{مثال} \\
\midrule
\code{int}    & \code{int}          & \code{128} \\
\code{float}  & \code{float}        & \code{0.1} \\
\code{bool}   & \code{bool}         & \code{true} / \code{false} \\
\code{string} & \code{str}          & \code{"relu"} \\
\code{shape}  & \code{tuple[int]}   & \code{(64, 32, 32)} \\
\code{None}   & \code{NoneType}     & \code{None} \\
\bottomrule
\end{tabular}
\end{center}

تمایز بین \code{int} و \code{float} در زمان کامپایل بررسی می‌شود. مثلاً
\code{Dropout(p=1)} خطا می‌دهد چون \code{p} انتظار \code{float} دارد نه \code{int}.

\section{انواع لایه}

\subsection{لایه‌های پارامتردار}

این لایه‌ها در \code{\_\_init\_\_} به صورت \code{self.<id> = nn.X(...)} تولید می‌شوند.

\begin{center}
\begin{tabular}{lll}
\toprule
\textbf{کلیدواژه DSL} & \textbf{کلاس PyTorch} & \textbf{پارامترهای اصلی} \\
\midrule
\code{Linear}        & \code{nn.Linear}             & \code{in\_features}, \code{out\_features}, \code{bias} \\
\code{Conv2d}        & \code{nn.Conv2d}             & \code{in\_ch}, \code{out\_ch}, \code{kernel}, \code{stride}, \code{padding} \\
\code{Conv1d}        & \code{nn.Conv1d}             & \code{in\_ch}, \code{out\_ch}, \code{kernel}, \code{stride} \\
\code{BatchNorm2d}   & \code{nn.BatchNorm2d}        & \code{num\_features} \\
\code{LayerNorm}     & \code{nn.LayerNorm}          & \code{normalized\_shape} \\
\code{MaxPool2d}     & \code{nn.MaxPool2d}          & \code{kernel}, \code{stride} \\
\code{AvgPool2d}     & \code{nn.AvgPool2d}          & \code{kernel}, \code{stride} \\
\code{Dropout}       & \code{nn.Dropout}            & \code{p} \\
\code{Flatten}       & \code{nn.Flatten}            & \code{start\_dim}, \code{end\_dim} \\
\code{Embedding}     & \code{nn.Embedding}          & \code{num\_embeddings}, \code{embedding\_dim} \\
\code{MultiHeadAttn} & \code{nn.MultiheadAttention} & \code{embed\_dim}, \code{num\_heads} \\
\code{LSTM}          & \code{nn.LSTM}               & \code{input\_size}, \code{hidden\_size}, \code{num\_layers} \\
\code{GRU}           & \code{nn.GRU}                & \code{input\_size}, \code{hidden\_size} \\
\bottomrule
\end{tabular}
\end{center}

\subsection{فعال‌سازی‌ها}

\code{ReLU}، \code{Sigmoid}، \code{Tanh}، \code{GELU}، \code{Softmax(dim)}،
\code{LeakyReLU(negative\_slope)}، \code{ELU(alpha)}.

\subsection{عملیات گراف (بدون ماژول)}

این انواع در \code{\_\_init\_\_} خطی تولید نمی‌کنند و فقط در \code{forward} ظاهر می‌شوند:

\begin{center}
\begin{tabular}{lll}
\toprule
\textbf{کلیدواژه} & \textbf{پارامترها} & \textbf{کد تولیدشده} \\
\midrule
\code{Add}      & ---                           & \code{out = a + b} \\
\code{Concat}   & \code{dim: int}               & \code{out = torch.cat([a, b], dim=d)} \\
\code{Residual} & ---                           & \code{out = main + shortcut} \\
\code{Split}    & \code{chunks: int, dim: int}  & \code{parts = torch.chunk(x, n, dim=d)} \\
\bottomrule
\end{tabular}
\end{center}

\section{نحو یال}

\begin{dfnbox}{یال ساده}{edge-simple}
\begin{latin}\code{edge src -> dst}\end{latin}

یک یال جهت‌دار از گره \code{src} به گره \code{dst} تعریف می‌کند.
\end{dfnbox}

\begin{dfnbox}{یال برچسب‌دار}{edge-labeled}
\begin{latin}\code{edge src -> dst [label="name"]}\end{latin}

یک یال با برچسب رشته‌ای تعریف می‌کند. برچسب‌ها دو نقش دارند:
\begin{itemize}
    \item \code{"shortcut"}: یال اتصال کوتاه به گره \code{Residual()} را مشخص می‌کند.
    \item \code{"residual\_*"}: اتصال کوتاه به گره‌ای با دو ورودی (الگوی \en{Transformer}).
\end{itemize}
\end{dfnbox}

\section{کامنت‌ها}

\begin{latin}
\begin{verbatim}
// کامنت تک‌خطی
/* کامنت
   چندخطی */
\end{verbatim}
\end{latin}

% ================================================================
\chapter{پیاده‌سازی}

\section{معماری کامپایلر}

\begin{dfnbox}{خط لوله کامپایل}{pipeline}
کامپایلر چهار مرحله دارد:
\begin{enumerate}
    \item \textbf{تجزیه} (\en{Parsing}): \code{FileStream} $\to$ \code{NNGraphLexer} $\to$
    \code{NNGraphParser.program()} $\to$ درخت تجزیه
    \item \textbf{تحلیل معنایی} (\en{Semantic Analysis}): \code{ParseTreeWalker} +
    \code{NNGraphCustomListener} $\to$ مدل گراف
    \item \textbf{تولید کد} (\en{Code Generation}): \code{CodeGenerator.generate()} $\to$
    رشته Python
    \item \textbf{خروجی}: نوشتن به فایل \code{.py} یا چاپ در \code{stdout}
\end{enumerate}
\end{dfnbox}

\subsection{مسئولیت فایل‌ها}

\begin{center}
\begin{tabular}{ll}
\toprule
\textbf{فایل} & \textbf{نقش} \\
\midrule
\code{NNGraph.g4}             & منبع حقیقت گرامر (۵۶ خط) \\
\code{gen/NNGraphLexer.py}    & تولیدشده توسط ANTLR4 --- نباید ویرایش شود \\
\code{gen/NNGraphParser.py}   & تولیدشده توسط ANTLR4 --- نباید ویرایش شود \\
\code{gen/NNGraphListener.py} & \en{listener} پایه با متدهای خالی --- نباید ویرایش شود \\
\code{custom\_listener.py}    & تحلیل معنایی (۲۹۳ خط) \\
\code{code\_generator.py}     & تولید کد PyTorch (۲۳۹ خط) \\
\code{main.py}                & رابط CLI و تابع \code{compile\_nng()} (۷۳ خط) \\
\bottomrule
\end{tabular}
\end{center}

\section{تحلیل‌گر معنایی}

\begin{dfnbox}{NNGraphCustomListener}{custom-listener}
کلاس \code{NNGraphCustomListener(NNGraphListener)} در \code{custom\_listener.py}
تمام منطق تحلیل معنایی را پیاده می‌کند. حالت داخلی:
\begin{itemize}[noitemsep]
    \item \code{self.nodes}: \code{dict[id, \{type, params, line\}]}
    \item \code{self.edges}: \code{list[\{src, dst, label, line\}]}
    \item \code{self.config}: \code{dict[key, value]}
\end{itemize}
\end{dfnbox}

\subsection{انتشار مقدار در درخت}

متدهای \code{exitShape\_expr} و \code{exitValue} مقادیر را روی خود context ذخیره
می‌کنند:

\begin{latin}
\begin{verbatim}
def exitShape_expr(self, ctx):
    ints = [int(ctx.INT_LITERAL(i).getText())
            for i in range(ctx.INT_LITERAL().__len__())]
    ctx.pvalue = tuple(ints)
    ctx.ptype  = tuple
\end{verbatim}
\end{latin}

این رویکرد تضمین می‌کند که هر context پس از خروج از آن، مقادیر فرزندانش را
در دسترس دارد---الگویی که در پروژه EVM مرجع نیز استفاده شده است.

\subsection{قوانین اعتبارسنجی}

هشت قانون اعتبارسنجی در دو فاز اجرا می‌شوند:

\textbf{فاز اول (در زمان پیمایش درخت):}
\begin{enumerate}
    \item \textbf{شناسه تکراری گره}: در \code{exitNode\_decl}
    \item \textbf{ارجاع نامعتبر در یال}: در \code{exitEdge\_decl}
    \item \textbf{نوع لایه ناشناس}: در \code{exitNode\_decl}
    \item \textbf{عدم تطابق نوع پارامتر}: در \code{exitNode\_decl}
\end{enumerate}

\textbf{فاز دوم (در \code{exitGraph\_block}، پس از تکمیل گراف):}
\begin{enumerate}
    \setcounter{enumi}{4}
    \item \textbf{گره خروجی اعلان‌نشده}: در \code{\_validate\_output\_declared}
    \item \textbf{گره고립} (هشدار): در \code{\_validate\_orphans}
    \item \textbf{آریتی Residual != 2}: در \code{\_validate\_residual\_arity}
    \item \textbf{تشخیص دور با Kahn}: در \code{\_validate\_no\_cycles}
    \item \textbf{قابلیت دسترسی BFS}: در \code{\_validate\_reachability}
\end{enumerate}

\subsection{قالب پیام خطا}

\begin{genbox}{قالب خطا}
\begin{latin}
\texttt{[line L:C] SemanticError: <message>}
\end{latin}

\vspace{0.5em}
که \code{L} شماره سطر (از ۱) و \code{C} ستون توکن (از ۰) است.
هشدارها فقط در \code{stderr} چاپ می‌شوند و کامپایل را متوقف نمی‌کنند.
\end{genbox}

\subsection{\_FakeCtx برای خطاهای پس از پیمایش}

خطاهایی که در \code{exitGraph\_block} رخ می‌دهند context زنده ندارند. یک کلاس داخلی
\code{\_FakeCtx} شماره سطر را از رکورد اعلان گره حمل می‌کند:

\begin{latin}
\begin{verbatim}
class _FakeCtx:
    class start:
        line   = info["line"]
        column = 0
\end{verbatim}
\end{latin}

\section{مولد کد}

\begin{dfnbox}{CodeGenerator}{codegen}
کلاس \code{CodeGenerator} در \code{code\_generator.py} چهار گام داخلی دارد:
\begin{enumerate}[noitemsep]
    \item \code{\_compute\_graph()}: ساخت لیست‌های مجاورت و مرتب‌سازی توپولوژیک
    \item \code{\_assign\_vars()}: تخصیص نام متغیر Python به هر گره
    \item \code{\_emit\_init()} + \code{\_emit\_forward()}: تولید خطوط کد
    \item \code{\_assemble()}: ترکیب در سورس Python نهایی
\end{enumerate}
\end{dfnbox}

\subsection{الگوریتم تخصیص متغیر}

\begin{dfnbox}{اصل اساسی}{var-invariant}
\code{var\_at[n] = 'x'} فقط وقتی که \code{out\_degree[predecessor] == 1}.

وقتی یک گره به چند مقصد فرستاده می‌شود (\code{out\_degree > 1})، تنسور آن
هنوز توسط شاخه‌های دیگر نیاز است. بازنویسی \code{x} در این نقطه آن را خراب
می‌کند. به جای آن، گره بعدی نام شناسه خودش را به عنوان متغیر دریافت می‌کند.
\end{dfnbox}

چهار قانون کامل (به ترتیب در \code{\_assign\_vars}):

\begin{latin}
\begin{verbatim}
1. گره ورودی:  var_at[input_id] = 'x'

2. گره ادغام (in_degree > 1):
       var_at[n] = 'x'       if out_degree[n] <= 1
       var_at[n] = node_id   if out_degree[n] > 1

3. شروع شاخه (predecessor.out_degree > 1):
       var_at[n] = node_id

4. ادامه خطی (predecessor.out_degree == 1):
       var_at[n] = var_at[predecessor]
\end{verbatim}
\end{latin}

\begin{exbox}{ردیابی MLP (زنجیره خطی)}{trace-mlp}
\begin{latin}
\texttt{x → fc1 → relu1 → fc2 → relu2 → drop → fc3 → out}
\end{latin}

هر predecessor دارای \code{out\_degree == 1} است $\Rightarrow$ قانون ۴ همیشه اعمال
می‌شود $\Rightarrow$ همه گره‌ها \code{var\_at = 'x'} دارند.
کد \code{forward} از \code{x} در همه جا استفاده می‌کند.
\end{exbox}

\begin{exbox}{ردیابی ResBlock (فن‌اوت سپس ادغام)}{trace-resblock}
\code{x} به \code{conv1} (مسیر اصلی) و \code{skip} (اتصال کوتاه) فرستاده می‌شود.
پس \code{out\_degree[x] = 2}.

\begin{itemize}
    \item \code{conv1}: قانون ۳ $\Rightarrow$ \code{var\_at = 'conv1'}
    \item \code{bn1}, \code{relu1}, \code{conv2}, \code{bn2}: قانون ۴ $\Rightarrow$ همه \code{'conv1'}
    \item \code{skip} (Residual، \code{in\_degree=2}): قانون ۲، \code{out\_degree=1}
    $\Rightarrow$ \code{var\_at = 'x'}
    \item کد تولیدشده: \code{x = conv1 + x}
\end{itemize}
\end{exbox}

\begin{exbox}{ردیابی InceptionBlock (دو شاخه موازی)}{trace-inception}
\code{out\_degree[x] = 2} (به \code{convA1} و \code{convB1}).

\begin{itemize}
    \item \code{convA1}: قانون ۳ $\Rightarrow$ \code{'convA1'}
    \item \code{convB1}: قانون ۳ $\Rightarrow$ \code{'convB1'}
    \item \code{convA2}: قانون ۴ $\Rightarrow$ ارث از \code{'convA1'}
    \item \code{convB2}: قانون ۴ $\Rightarrow$ ارث از \code{'convB1'}
    \item \code{cat1} (Concat، \code{in\_degree=2}): قانون ۲ $\Rightarrow$ \code{'x'}
    \item کد: \code{x = torch.cat([convA1, convB1], dim=1)}
\end{itemize}
\end{exbox}

\begin{exbox}{ردیابی TransformerEncoder (دو زیرلایه residual ضمنی)}{trace-transformer}
\code{out\_degree[x] = 2} (به \code{attn} و \code{norm1}).\\
\code{out\_degree[norm1] = 2} (به \code{ff1} و \code{norm2}).

\begin{itemize}
    \item \code{attn}: قانون ۳ $\Rightarrow$ \code{'attn'}
    \item \code{drop1}: قانون ۴ $\Rightarrow$ \code{'attn'}
    \item \code{norm1} (\code{in\_degree=2, out\_degree=2}): قانون ۲ $\Rightarrow$ \code{'norm1'}
    \item \code{ff1}: قانون ۳ $\Rightarrow$ \code{'ff1'}
    \item \code{gelu1, drop2, ff2}: قانون ۴ $\Rightarrow$ \code{'ff1'}
    \item \code{norm2} (\code{in\_degree=2, out\_degree=1}): قانون ۲ $\Rightarrow$ \code{'x'}
\end{itemize}

خطوط تولیدشده:
\begin{latin}
\begin{verbatim}
attn, _ = self.attn(x, x, x)
attn     = self.drop1(attn)
norm1    = self.norm1(x + attn)
ff1      = self.ff1(norm1)
ff1      = self.gelu1(ff1)
ff1      = self.drop2(ff1)
ff1      = self.ff2(ff1)
x        = self.norm2(norm1 + ff1)
\end{verbatim}
\end{latin}
\end{exbox}

\subsection{موارد خاص تولید کد}

\begin{center}
\begin{tabular}{lll}
\toprule
\textbf{الگو} & \textbf{شرط} & \textbf{کد تولیدشده} \\
\midrule
\code{MultiHeadAttn}    & نوع لایه      & \code{out, \_ = self.attn(x, x, x)} \\
\code{LSTM / GRU}       & نوع لایه      & \code{out, \_ = self.lstm(x)} \\
\code{Residual}         & گره بدون ماژول & \code{out = main + shortcut} \\
\code{Add}              & گره بدون ماژول & \code{out = a + b} \\
\code{Concat}           & گره بدون ماژول & \code{out = torch.cat([a, b], dim=d)} \\
\code{Split}            & گره بدون ماژول & \code{parts = torch.chunk(x, n, dim=d)} \\
\en{Residual ضمنی}   & \code{in\_degree>1} + برچسب \code{residual\_*} & \code{out = self.norm(skip + main)} \\
\bottomrule
\end{tabular}
\end{center}

\begin{notebox}
\textbf{نگاشت نام پارامتر:} DSL از نام‌های کوتاه‌تر استفاده می‌کند. مولد کد آن‌ها
را به نام‌های PyTorch نگاشت می‌کند: \code{in\_ch} $\to$ \code{in\_channels}،
\code{out\_ch} $\to$ \code{out\_channels}، \code{kernel} $\to$ \code{kernel\_size}.
همه نام‌های دیگر بدون تغییر عبور می‌کنند.
\end{notebox}

\subsection{batch\_first در MultiHeadAttn}

\code{nn.MultiheadAttention} از \en{PyTorch} به صورت پیش‌فرض \code{batch\_first=False}
است. مولد کد همیشه \code{batch\_first=True} اضافه می‌کند تا با شکل رایج
\code{(batch, seq, dim)} سازگار باشد.

\subsection{ساختار کد خروجی}

\begin{latin}
\begin{verbatim}
import torch
import torch.nn as nn


class <ModelName>(nn.Module):
    def __init__(self):
        super(<ModelName>, self).__init__()
        <init_lines>

    def forward(self, x):
        <forward_lines>
        return <output_var>


if __name__ == '__main__':
    device = torch.device('<device>')
    model  = <ModelName>().to(device)
    x      = torch.randn(<batch_size>, <input_shape...>).to(device)
    print(model(x).shape)
\end{verbatim}
\end{latin}

% ================================================================
\chapter{ارزیابی}

\section{مجموعه آزمایش}

\begin{dfnbox}{آمار آزمایش}{test-stats}
\begin{itemize}[noitemsep]
    \item تعداد کل تست‌ها: \textbf{۵۵}
    \item فریم‌ورک: \code{pytest} نسخه ۸.۴.۲
    \item زمان اجرا: حدود ۱.۰۹ ثانیه
\end{itemize}
\end{dfnbox}

\begin{center}
\begin{tabular}{lcp{6cm}}
\toprule
\textbf{فایل} & \textbf{تست‌ها} & \textbf{پوشش} \\
\midrule
\code{test\_parser.py}   & ۸  & گرامر بدون خطای نحوی برای ۴ نمونه؛ config، برچسب، شکل \\
\code{test\_listener.py} & ۱۴ & ۴ ورودی معتبر؛ ۱۰ شرط خطا با \code{SystemExit} \\
\code{test\_codegen.py}  & ۱۹ & نام کلاس‌های nn، مقادیر پارامتر، الگوهای forward \\
\code{test\_e2e.py}      & ۱۴ & کامپایل کامل + ایجاد نمونه + forward pass با شکل صحیح \\
\bottomrule
\end{tabular}
\end{center}

\section{نمونه‌های کارکرده}

\begin{exbox}{MLP — زنجیره خطی ۳ لایه}{eval-mlp}
\begin{itemize}[noitemsep]
    \item ورودی: \code{tensor(784)} --- خروجی: \code{(batch, 10)}
    \item توپولوژی: زنجیره خطی
    \item استراتژی متغیر: همه \code{x}
    \item آزمایش: \code{Softmax} ردیف‌ها به ۱ جمع می‌شوند (\code{atol=1e-5})
\end{itemize}
\end{exbox}

\begin{exbox}{ResBlock — بلوک کانولوشنی با اتصال باقیمانده}{eval-resblock}
\begin{itemize}[noitemsep]
    \item ورودی: \code{tensor(64, 32, 32)} --- خروجی: \code{(batch, 65536)} پس از Flatten
    \item توپولوژی: فن‌اوت از ورودی، ادغام در \code{Residual()}
    \item آزمایش: شکل صحیح؛ در حالت eval خروجی قطعی است
\end{itemize}
\end{exbox}

\begin{exbox}{InceptionBlock — دو شاخه Conv موازی}{eval-inception}
\begin{itemize}[noitemsep]
    \item ورودی: \code{tensor(32, 56, 56)} --- خروجی: \code{(batch, 48, 56, 56)}
    \item شاخه A: ۱۶ کانال؛ شاخه B: ۳۲ کانال؛ ادغام: \code{Concat(dim=1)} $\to$ ۴۸ کانال
    \item آزمایش: ابعاد فضایی تغییر نمی‌کنند
\end{itemize}
\end{exbox}

\begin{exbox}{TransformerEncoder — یک لایه encoder}{eval-transformer}
\begin{itemize}[noitemsep]
    \item ورودی: \code{tensor(512, 128)} --- خروجی: \code{(1024, 128)} پس از Flatten
    \item دو زیرلایه residual ضمنی (بدون گره \code{Residual()})
    \item \code{MultiHeadAttn} با \code{embed\_dim=128, num\_heads=8}
    \item آزمایش: شکل صحیح؛ در حالت eval خروجی قطعی است
\end{itemize}
\end{exbox}

\section{اثبات درستی متغیرها}

\begin{dfnbox}{اثبات عدم بازنویسی تنسور}{var-correctness}
برای هر نمونه شاخه‌دار، \code{data\_ptr()} تنسور ورودی در طول محاسبات forward
ردیابی شد:
\begin{itemize}[noitemsep]
    \item \textbf{ResBlock}: \code{x.data\_ptr()} در نقطه \code{x = conv1 + x} تغییر نکرده
    \item \textbf{Inception}: \code{x.data\_ptr()} هنگام فراخوانی \code{convA1(x)} و \code{convB1(x)} یکسان است
    \item \textbf{Transformer}: \code{x.data\_ptr()} هنگام \code{norm1(x + attn)}؛ \code{norm1.data\_ptr()} هنگام \code{norm2(norm1 + ff1)} تغییر نکرده
\end{itemize}
\end{dfnbox}

این نتیجه مستقیماً از اصل \code{var\_at[n]='x'} فقط وقتی \code{out\_degree=1} پیروی می‌کند.

\section{آمار کدبیس}

\begin{center}
\begin{tabular}{lr}
\toprule
\textbf{فایل/دایرکتوری} & \textbf{خطوط} \\
\midrule
\code{NNGraph.g4}                    & ۵۶ \\
\code{custom\_listener.py}           & ۲۹۳ \\
\code{code\_generator.py}            & ۲۳۹ \\
\code{main.py}                       & ۷۳ \\
\code{tests/} (۴ فایل)              & ۵۷۵ \\
\code{required\_code\_collection/}   & $\approx$۱۰۰ \\
\midrule
\textbf{جمع}                         & $\approx$۱۳۳۶ \\
\bottomrule
\end{tabular}
\end{center}

\begin{itemize}[noitemsep]
    \item زبان: Python 3.14.3
    \item ANTLR: 4.13.2
    \item PyTorch: 2.10.0
    \item pytest: 8.4.2
\end{itemize}

% ================================================================
\chapter{تصمیمات طراحی}

\section{Listener در برابر Visitor}

الگوی Listener به صورت طبیعی با انباشت حالت سازگار است. متدهای \code{exit*}
می‌توانند مستقیماً روی \code{self.nodes} و \code{self.edges} عمل کنند بدون اینکه
مقادیر را از طریق درخت برگردانند. برای یک کامپایلر گراف که هدفش ساختن یک مدل
گراف از درخت تجزیه است، این انتخاب طبیعی‌تر از Visitor است.

\section{load\_from\_listener() در برابر generate(traversal)}

پروژه EVM مرجع از \code{CodeGenerator.generate(traversal)} استفاده می‌کند که یک
لیست توکن پس‌ترتیب دریافت می‌کند---مناسب برای تولید بایت‌کد پشته‌ای. برای
کامپایلر گراف، ورودی طبیعی ساختار گراف (گره‌ها + یال‌ها) است نه جریان خطی
توکن. پس \code{load\_from\_listener(listener)} مستقیم‌تر است.

\section{نام‌گذاری متغیر با node\_id}

مثال‌های مشخصات زبان از نام‌های تک‌حرفی (\code{a}، \code{b}) یا نام‌های توصیفی
(\code{identity}، \code{attn\_out}) استفاده می‌کنند. این کامپایلر از شناسه واقعی گره
(\code{convA1}، \code{conv1}، \code{attn}) استفاده می‌کند زیرا:
\begin{itemize}[noitemsep]
    \item قطعی است
    \item خوانا است
    \item برای گراف‌های بزرگ‌تر بدون شمارنده کار می‌کند
    \item خروجی ریاضی یکسان است
\end{itemize}

% ================================================================
\chapter{نتیجه‌گیری و کارهای آینده}

\section{نتیجه‌گیری}

\en{NNGraph DSL} یک زنجیره کامپایلر چهار مرحله‌ای از \code{.nng} به
\code{nn.Module} قابل اجرا ارائه می‌دهد. کامپایلر:

\begin{itemize}
    \item گرامر ANTLR4 با ۱۴ قانون تجزیه‌گر و ۲۰ توکن
    \item ۲۴ نوع لایه را پشتیبانی می‌کند
    \item ۸ قانون اعتبارسنجی را اجرا می‌کند
    \item شاخه‌بندی پیچیده را بدون بازنویسی تنسور مدیریت می‌کند
    \item با ۵۵ تست در چهار معماری نمونه اعتبارسنجی شده است
\end{itemize}

\section{کارهای آینده}

این موارد در مشخصات زبان پیش‌بینی شده‌اند اما پیاده‌سازی نشده‌اند:

\begin{center}
\begin{tabular}{ll}
\toprule
\textbf{قابلیت} & \textbf{توضیح} \\
\midrule
استنتاج شکل        & انتشار شکل تنسور در طول گراف و تشخیص عدم تطابق در کامپایل \\
صادرکردن ONNX       & تولید گراف سازگار با ONNX به جای Python \\
داربست آموزش        & تولید حلقه آموزش با بهینه‌ساز و تابع خطا از بلوک \code{config} \\
تجسم Graphviz       & خروجی نمودار گراف مدل \\
زیرگراف‌های نام‌دار  & بلوک‌های ماکرو قابل استفاده مجدد (مثلاً \code{ResBlock} قابل تکرار) \\
کنترل جریان پویا   & شاخه‌های \code{if}/\code{else} درون \code{forward} \\
حاشیه‌نویسی کوانتیزاسیون & متادیتا برای تبدیل مدل به نسخه‌های کوانتیزه \\
\bottomrule
\end{tabular}
\end{center}

% ================================================================
\chapter*{منابع}
\markboth{منابع}{منابع}
\addcontentsline{toc}{chapter}{منابع}

\begin{enumerate}
    \item \en{Parr, T. (2013). \textit{The Definitive ANTLR 4 Reference}. Pragmatic Bookshelf.}
    \item \en{Paszke, A., et al. (2019). PyTorch: An Imperative Style, High-Performance
    Deep Learning Library. \textit{NeurIPS 2019}.}
    \item \en{NADER Framework: Neural Architecture Design via Representation (referenced
    in \texttt{NNGraph\_DSL\_Specification.md}).}
    \item \en{Zhang, A. M. \textit{amznotes LaTeX template}.
    \texttt{github.com/alexmingzhang/latex-notes-template}}
\end{enumerate}

\end{document}
